% =====================================================================
%  Secao 5: Soluções Parciais
% =====================================================================
\section{Soluções Parciais}

Uma das ideias mais importantes do Backtracking é o conceito de solução parcial.

Considere o objetivo de montar uma sequência formada por quatro números. Ao invés de tentar construir a sequência completa imediatamente, o algoritmo trabalha em etapas.

Por exemplo:

\begin{passos}[Construção passo a passo]
\begin{description}[leftmargin=4.2em,font=\normalfont\bfseries]
  \item[Etapa 1:] $(3)$
  \item[Etapa 2:] $(3,\ 1)$
  \item[Etapa 3:] $(3,\ 1,\ 4)$
  \item[Etapa 4:] $(3,\ 1,\ 4,\ 2)$
\end{description}
\end{passos}

Em cada etapa o algoritmo verifica se ainda vale a pena continuar explorando aquele caminho. Caso identifique um problema, retorna para a etapa anterior.

\figura{fig:solucao-parcial}%
  {O Backtracking constrói a solução gradualmente, validando cada nova escolha antes de prosseguir.}%
  {figuras/tikz/solucao-parcial}
